\documentclass[11pt,a4paper]{article}
\usepackage[margin=24mm]{geometry}
\usepackage{amsmath,amssymb,amsthm,booktabs,array}
\usepackage[hidelinks]{hyperref}
\newcommand{\pos}[1]{[#1]_+}
\newcommand{\dd}{\,\mathrm d}
\newcommand{\kc}{k_c}
\newcommand{\bud}{\mathcal B}
\newtheorem{proposition}{Proposition}
\title{AS1: independent initial lane loads\newline and the existing forward model}
\author{Elementary VSL lane-drop task review}
\date{19 October 2026}
\begin{document}
\maketitle
\noindent\textbf{Status.} Bounded AS1 review completed; exact amendment
submitted at its user checkpoint. Step 3 remains open. AS2's observation
contract and AS3's action/gradient conclusions are not established here.

\paragraph{Traffic meaning.}
The same ending-lane population can face a crowded adjacent lane or an
uncrowded adjacent lane even when the continuing total is identical.
Allowing that distinction requires independent initial lane loads, not a
different traffic mechanism. The existing H1 preparation, finite T1
terminal, S-E-P allocation, R-E source, O-R renewal and W-F work laws
support the asymmetric family below under their retained sufficient bounds.
This is a conditional mathematical qualification, not empirical admission.

\section{The exact initialization amendment}
E is the ending lane, A its adjacent continuing lane and B the far continuing
lane. The same uniform longitudinal footprint $[0,\ell]$ occupies the
preparation reach at the decision, with $0<\ell\leq L_p$; the terminal of
length $b>0$ and the recovery account are empty. Set $L=L_p+b$ and
\begin{equation}
 e_0,a_0,b_0\geq0,\qquad r=e_0+a_0,\qquad
 \rho_0=e_0+a_0+b_0=r+b_0,\qquad M=\ell\rho_0>0.
 \label{eq:initial}
\end{equation}
There is \emph{no} requirement $a_0=b_0$. Equal loading remains an admitted
special case. Longitudinal homogeneity within each lane is retained.
There are no postdecision external arrivals. One common constant advice
$v\in[v_-,v_+]$, $0<v_-<v_+$, is applied to each fixed true initialization.

Define the same fundamental diagram and inverse boundary mark:
\begin{equation}
 \begin{gathered}
 K=wk_j,\quad \kc(v)=\frac{K}{w+v},\quad C(v)=\frac{Kv}{w+v},\\
 c(u)=\frac{Ku}{w+u},\qquad U(g)=\frac{wg}{K-g}.
 \end{gathered}
 \label{eq:fd}
\end{equation}
Retain all existing sufficient restrictions, now with the actual three-lane
mass in \eqref{eq:initial}:
\begin{align}
 &w,k_j,h,a_r>0,\quad\gamma\geq1,\quad
 e_0,a_0,b_0\leq\kc(v_+),\nonumber\\
 &v_+(e_0+a_0)<R\leq C(v_-),\quad
 v_+b_0\leq R_B\leq C(v_-),\quad M<k_jb.
 \label{eq:domain}
\end{align}
The fixed physical coefficient and root condition are unchanged:
\begin{equation}
 \begin{gathered}
 c_0=1+\delta,\quad \delta,\beta\geq0,\quad
 \chi=\frac{\beta}{C_\star v_\star},\quad
 v_\star=v_+,\quad C_\star=C(v_\star),\\
 \eta=c_0-\beta(1+v_+/w)>0,\qquad c_{\max}=c_0+\beta.
 \end{gathered} \label{eq:workdomain}
\end{equation}
The anchor defining $\chi$ stays fixed across advice; $m=1$ veh-eq is the
same physical mass normalization. All geometry and physical coefficients
are shared across the contrasting initial states. The initial-dependent
$r$ is derived true state, not automatically a shared observer input.
Density has units veh-eq/length, rates have units veh-eq/time, and $M$ has
units veh-eq; $b_0$ is a density and $b$ a length.

\section{Preparation, handoff and storage: no lane equality is used}
The retained rule is
\begin{equation}
 \begin{gathered}
 S(v,a)=\min\{C(v),w(k_j-a)\},\qquad
 f=\min\{ve,(S(v,a)-va)/\gamma\},\\
 \dot e=-f/h,\quad\dot a=f/h,\quad\dot b_0=0.
 \end{gathered}
 \label{eq:h1}
\end{equation}
The transfer cancels in $e+a=r$. For every admitted advice,
\begin{equation}
 r<R/v_+\leq R/v\leq C(v)/v=\kc(v).
 \label{eq:free}
\end{equation}
Consequently $0\leq e\leq e_0$, $a=r-e<\kc(v)$ and $S=C(v)$.
The scalar right-hand side is Lipschitz and nonnegative toward the invariant
boundaries. No equality with B is needed. With progress $z=vt$ and
$\Delta(v)=\kc(v)-r>0$,
\begin{equation}
 e_z=-\frac1h\min\{e,(\Delta+e)/\gamma\},\qquad a=r-e.
 \label{eq:progress}
\end{equation}
The two explicit branches remain
\begin{equation}
 e=e_s e^{-(z-z_s)/h}\quad\hbox{or}\quad
 e=(e_s+\Delta)e^{-(z-z_s)/(\gamma h)}-\Delta,
 \label{eq:branches}
\end{equation}
with a receiving-to-donor switch at $e=\Delta/(\gamma-1)$ for
$\gamma>1$. For $\gamma=1$ the donor branch is active throughout.

Let $t_f^p=(L_p-\ell)/v$, $t_r^p=L_p/v$. During that pulse only,
\begin{equation}
 \dot F_E=ve(t),\quad \dot F_A=v(r-e(t)),\quad\dot F_B=vb_0.
 \label{eq:entry}
\end{equation}
Let $s(t)=\ell$ before $t_f^p$, $s=L_p-vt$ during the pulse, and $s=0$
after $t_r^p$. With $J_{\rm lat}=\int_0^t s f/h\dd u$, the balances are
\begin{align}
 se+F_E+J_{\rm lat}&=\ell e_0,&
 sa+F_A-J_{\rm lat}&=\ell a_0,\nonumber\\
 sb_0+F_B&=\ell b_0,&s\rho_0+F_E+F_A+F_B&=M.
 \label{eq:massprep}
\end{align}
Each identity follows from its initial value and its derivative before,
during and after the pulse. Early-merged mass enters A once and is not later
counted as forced E service.

For $j\in\{E,A,B\}$, retain
\begin{equation}
 P_j(t)=F_j(t-b/v),\quad N_j=F_j-Y_j,\quad q_j=P_j-Y_j,
 \quad \lambda_j=\dot P_j,\quad g_j=\dot Y_j,
 \label{eq:counts}
\end{equation}
with zero negative-time histories. $N_j$ is physical terminal stock;
$q_j$ is eligible waiting mass. The far lane has $Y_B=P_B$ since
$vb_0\leq R_B$. Its density is not derived from A.

The terminal entrance vacancy is
\begin{equation}
 V_j^{\rm vac}(t)=
 \begin{cases}k_jwt-F_j(t),&0<t\leq b/w,\\
 Y_j(t-b/w)+k_jb-F_j(t),&t>b/w.
 \end{cases} \label{eq:vacancy}
\end{equation}
All entry rates are at most $C(v)<wk_j$. The early expression is positive;
the later expression is at least $k_jb-M>0$. Thus finite receiving accepts
the preparation handoff under the retained conservative total-mass bound.
This protects the homogeneous upstream reduction; it is not a proof of
general spillback dynamics outside this sufficient domain.

The eligible pulse spans $t_f=(L-\ell)/v$ to $t_r=L/v$. Inside it,
\begin{equation}
 \lambda_A+\lambda_E=vr<R,\qquad
 \lambda'_A=-\lambda'_E\in[0,L_A],\qquad L_A=v_+^2e_0/h.
 \label{eq:pulse}
\end{equation}
E offers decrease and A offers increase. These facts use $e+a=r$ and the
donor upper bound $f\leq ve$, not $a_0=b_0$.

\section{The selected point service and recovery are unchanged}
Free marks are $u_j=v$ and queued marks $u_j=U(g_j)$. Retain
\begin{equation}
 c_E=c_0+\chi g_A\pos{u_A-u_E},\qquad
 H(g;a,u)=g\{c_0+\chi a\pos{u-U(g)}\}.
 \label{eq:root}
\end{equation}
Here the argument $a$ in $H$ denotes an A \emph{flow}, not the upstream
density. On the selected feasible interval, $H_g\geq\eta>0$ and
$x/c_{\max}\leq H^{-1}(x;a,u)\leq x/c_0$. The proof uses bounds on
flows, marks and fixed coefficients, not initial lane equality.

The S-E-P branches are retained exactly:
\begin{center}\small
\begin{tabular}{p{28mm}p{118mm}}
\toprule
Current mode & Service and free-feasibility condition\\
\midrule
A/E free & $g_A=\lambda_A$, $g_E=\lambda_E$ if
$\lambda_A+c_0\lambda_E\leq\bud$.\\[4pt]
A free, E queued & $g_A=\lambda_A\leq\bud$,
$H(g_E;\lambda_A,v)=\bud-\lambda_A$.\\[4pt]
A/E queued & $g_A=\bud/2$,
$H(g_E;\bud/2,U(\bud/2))=\bud/2$.\\[4pt]
A queued, E free & $g_E=\lambda_E$, $g_A=\bud-c_0\lambda_E$
if $c_0\lambda_E\leq\bud/2$; otherwise E queues.\\
\bottomrule
\end{tabular}\end{center}
Retain strict-shortfall queue birth, actual-zero free-admission retesting,
maximal feasible free continuation at equality, and right-hand data at
arrival jumps. There is no new priority, timer, smoothing or branch latch.

For E's insertion mark $u$, retain
\begin{align}
 V(\theta;u)&=\sqrt{(w+u)^2+2wa_r\theta}-w,\nonumber\\
 D(\theta;u)&=\pos{R-c(V(\theta;u))},\nonumber\\
 p(t)&=\frac1m\int_{\pos{Y_E(t)-m}}^{Y_E(t)}
 D(t-T(\xi);u_E(T(\xi)))\dd\xi,\qquad \bud=\pos{R-p}.
 \label{eq:recovery}
\end{align}
$T$ is E's mass-rank admission-time inverse. Only E deposits loss. Recent
mass replaces oldest retained weight; profiles also age. The account
retains history, not just its current reading, and it does not keep the
strictest prior cap. $V$ is an equivalent boundary-speed quantity, not a
vehicle trajectory. Neither the overlap approximation nor its physical
validity is strengthened by this initialization amendment.

\section{Forward qualification carries to every admitted asymmetric state}
\begin{proposition}
For \eqref{eq:initial}--\eqref{eq:workdomain}, the unchanged laws and causal
guards have the same conditional unique complete forward continuation for
every advice in the original interval. Mass, feasible service and finite
terminal storage are respected. All traffic clears in finite time.
\end{proposition}
\begin{proof}
The preceding preparation/storage argument supplies the same bounded,
monotone eligible pulse as the accepted W-F proof. To make the remaining
dependencies explicit, put $G=R/c_0$ and retain its bounds
\begin{equation}
 p\leq p_{\max}:=\frac Gm\int_0^{m/G}D(\theta;0)\dd\theta<R,
 \quad B_{\min}=R-p_{\max}>0,\quad
 |\dot{\bud}|\leq L_B:=2GR/m+a_rk_j.
 \label{eq:budgetbounds}
\end{equation}
The account bound fills the youngest possible ages at the maximal admitted
E rate. Strictness uses $a_r>0$, not an imposed service floor. The history
comparison for identical past and two subsequent flow/mark histories is
\begin{equation}
 |\Delta p(t)|\leq\frac1m\left(5R\int_{t_0}^t|\Delta g_E|\dd s
                 +k_jG\int_{t_0}^t|\Delta u_E|\dd s\right).
 \label{eq:volterra}
\end{equation}
Together with $H_g\geq\eta$, this gives the same short-interval contraction
on each valid mode. It does not assume a positive old E rate or initial
symmetry. No numerical timestep is being defined.

At the front A's ordinary offer is below $R$. If joint free work fits, free
E has $v\geq U(R)$ and produces no recovery loss. Joint work
$vr+\delta\lambda_E$ decreases, so both remain free. Otherwise E's strict
shortfall gives a compatible queue birth. At any later E clearance with
free A, $\lambda_A+c_0\lambda_E\leq\bud$; with queued A,
$c_0\lambda_E\leq\bud/2$. Free E then adds no loss, old loss ages/retires,
and its offer decreases. These tests persist through any A clearance.
After the rear there are no arrivals. E has at most one waiting episode
under these specific assumptions, not under finite-cohort initialization
alone.

While E queues, A birth is at $\lambda_A=\bud$, with right-hand queue
growth $\bud/2\geq B_{\min}/2$. A clearance requires
$\lambda_A\leq\bud/2$, leaving a free margin at least $B_{\min}/2$.
The bounds \eqref{eq:pulse} and \eqref{eq:budgetbounds} give a common
positive separation for repeated A births/clearances,
\begin{equation}
 \tau_0=\frac{B_{\min}}{2(L_B+2L_A)}>0.
 \label{eq:separation}
\end{equation}
E clearance and the front/rear add only finitely many exceptions.
The same feasible guards therefore join the unique local histories without
event accumulation. After the rear, whenever $Q=q_A+q_E>0$,
$g_A+g_E\geq B_{\min}/c_{\max}$. Hence
\begin{equation}
 T_{\rm clear}(v)\leq L/v+c_{\max}M/B_{\min},\quad
 T_{\max}=L/v_-+c_{\max}M/B_{\min}
 \label{eq:clearance}
\end{equation}
is a common sufficient horizon. Every inequality depends on $r$, separate
B admission, total mass or common coefficients. None requires $a_0=b_0$.
This transfers the accepted conditional forward argument; it does not
establish smooth advice dependence at all queue boundaries.
\end{proof}

The complete balances and local total system travel time are consequently
\begin{align}
 s\rho_0+\sum_jN_j+\sum_jY_j&=M,\nonumber\\
 J(v)=\int_0^{T_{\max}}\left(M-\sum_jY_j(t)\right)\dd t
 &=\frac{M(L-\ell/2)}v+\int_0^\infty(q_A+q_E)\dd t.
 \label{eq:tst}
\end{align}
The first term integrates each initial position's remaining distance
$(L-x)$ over all three lanes. It needs $\rho_0=e_0+a_0+b_0$, not symmetry.
B has no eligible queue in this domain. Completion remains at the choke;
retained effects of completed E traffic affect unfinished followers only.

\section{A nonempty fixed-E, fixed-continuing-total asymmetric family}
In general, hold $e_0=e$ and $a_0+b_0=c$ fixed and set
\begin{equation}
 a_0=c/2+d,\qquad b_0=c/2-d,\qquad r=e+c/2+d,
 \quad M=\ell(e+c).
 \label{eq:split}
\end{equation}
The allowed $d$ obeys the existing bounds
\begin{equation}
 \begin{gathered}
 -c/2\leq d\leq c/2,\quad
 c/2+d\leq\kc(v_+),\quad c/2-d\leq\kc(v_+),\\
 v_+(e+c/2+d)<R,\qquad v_+(c/2-d)\leq R_B,
 \end{gathered} \label{eq:splitdomain}
\end{equation}
along with $e\leq\kc(v_+)$ and $\ell(e+c)<k_jb$ and the common parameter
bounds. This is not permission to use arbitrary imbalance at any load.

For an exact existence example, choose arbitrary positive speed scale $u$,
density scale $\kappa$, lengths $\ell,h$, and acceleration $a_r$. Put
\begin{equation}
 \begin{gathered}
 v_-=u,\ v_+=2u,\ w=2u,\ k_j=3\kappa,\ \gamma=2,
 \quad R=39\kappa u/20,\quad R_B=\kappa u,\\
 L_p=2\ell,\quad b=\ell,\quad
 c_0=11/10,\quad\beta=1/100,\quad
 e_0=3\kappa/5,\quad c=11\kappa/20,\\
 a_0=11\kappa/40+d,\quad b_0=11\kappa/40-d,
 \qquad |d|\leq3\kappa/40.
 \end{gathered} \label{eq:example}
\end{equation}
Then $r\in[4\kappa/5,19\kappa/20]$,
$M=23\kappa\ell/20$, $\kc(v_+)=3\kappa/2$ and $C(v_-)=2\kappa u$.
The worst ordinary E/A margin is
$R-v_+r\geq\kappa u/20>0$; the far margin is
$R_B-v_+b_0\geq3\kappa u/10>0$;
$M/(k_jb)=23/60<1$; and $\eta=27/25>0$.
All initial lanes have positive free-flow density with substantial margins.

At the two endpoints the A/B loads are swapped. At $v=v_+$ and the common
initial preparation progress, the two alternatives in \eqref{eq:h1} are:
\begin{center}
\begin{tabular}{llll}
\toprule
$(a_0,b_0)/\kappa$ & Donor offer $ve_0$ & Receiving alternative & Selected $f$\\
\midrule
$(1/5,7/20)$ & $6\kappa u/5$ & $13\kappa u/10$ & $6\kappa u/5$\\
$(7/20,1/5)$ & $6\kappa u/5$ & $23\kappa u/20$ & $23\kappa u/20$\\
\bottomrule
\end{tabular}\end{center}
The first preparation is donor-limited and the second receiving-limited.
The inequalities are strict, so this rate contrast persists over a positive
initial progress interval by continuity. The actual density source is
$f/h$. This is exact local algebra, with no traffic rollout or calibrated
parameter recommendation. It establishes a genuine mechanistic distinction
inside the admitted family, not indistinguishable histories, different
optimal advice or useful gradients.

\section{What transfers, what stops here, and provenance}
The whole-forward result, complete mass accounting, own-lane-first
preparation, monotone pulse, storage protection, roots, recent-mass recovery
and TST transfer as stated above. The old initial-E-only reconstruction
does not: E and $M$ fix $a_0+b_0$, not $a_0$. Symmetric matched-history and
action examples remain valid within their old member but cannot certify
the new Level 2-to-Level 3 target. AS2 must examine every permitted
observation and history, including lawful inference. AS3 must separately
qualify action relevance and the complete advice derivative. Neither an
interior optimum nor a positive information advantage is guaranteed.

\paragraph{Reused derivations.}
This is an original dependency audit and initialization extension of the
accepted review model. The supporting lineage is
\href{vsl_feedback8_full_mathematics.pdf}{the feedback-8 consolidated mathematics},
Sections 2--11 (initialization through global forward/TST), and
\href{vsl_wf_forward_qualification.pdf}{the W-F whole-domain qualification}.
The precise source locations are retained separately in the
\href{../../extraction/research_readiness/vsl_as1_initialization_proof_map.md}{AS1 component and proof dependency map}.
No new paper was viewed or new literature-backed mechanism selected for
AS1. Existing literature supports the ingredients to the extent recorded
in its source audits; it is not attributed this complete asymmetric law or
the illustrative parameter values. Applied MP6 remains open. SUMO is
optional. This document ends at the AS1 contract-confirmation checkpoint.
\end{document}
